\section{Full Proof of Theorem~\ref{thm:regret}}
\label{app:proof}

The proof proceeds by a four-lemma layering. \textbf{L1} derives a
two-resource empirical-Bernstein concentration inequality that is used
uniformly across all arms and resources. \textbf{L2} shows that
shrinking the budget to $B' = B - z\,\sigma\sqrt{T}$ preserves primal
feasibility with probability at least $1-\delta$. \textbf{L3} reduces
the chance-constrained primal to an unconstrained \BwK-style
primal--dual problem on a \emph{modified} instance whose utility
function is the per-unit-resource Lagrangian
$\mathcal{L}_i(\lambda)$. \textbf{L4} applies the \BwK~regret bound to
the modified instance and translates it back to the original objective.
Constants are kept explicit throughout so the composition can be
audited.

\paragraph{Notation and setup.}
Let $N_i(t)$ be the number of pulls of arm $i$ up to round $t$. Write
$\hat{c}_i(t), \hat{\ell}_i(t), \hat{u}_i(t)$ for the empirical means,
and $\sigma^2_c, \sigma^2_\ell, \sigma^2_u$ for the sub-Gaussian proxy
variances. Let $V_i(t)$ denote the empirical variance of cost
observations on arm $i$ up to round $t$.

\paragraph{L1 (two-resource empirical-Bernstein concentration).}
Under assumption (A1), for any $\delta' \in (0, 1)$ and any arm $i$
with $N_i \ge 2$,
\begin{equation}
\label{eq:L1}
  \Prob\!\bigl[|\hat{c}_i - \Ex c_i| > r_i(\delta')\bigr] \le \delta',
  \qquad
  r_i(\delta')
   \;=\;
    \sqrt{\tfrac{2 V_i \log(2/\delta')}{N_i}}
     + \tfrac{7 M \log(2/\delta')}{3 (N_i - 1)},
\end{equation}
and analogously for $\hat{\ell}_i, \hat{u}_i$.
\emph{Proof.} Eq.~\eqref{eq:L1} is the empirical-Bernstein bound
(Maurer--Pontil, Audibert--Munos--Szepesv\'ari). A union bound over the
three resources and $K$ arms gives a uniform bound at level $3K\delta'$;
setting $\delta' = \delta/(3KT)$ and union-bounding over $T$ rounds
yields a uniform confidence radius
$r(t) = \sqrt{2 V_i \log(6KT/\delta) / N_i(t)}
      + 7M\log(6KT/\delta)/(3(N_i(t)-1))
      = \Oh(\sqrt{\log(T/\delta)/N_i(t)})$
that holds with probability at least $1-\delta/3$. \qed

\paragraph{L2 (feasibility preservation under budget shrinkage).}
Let $x^\star \in \{0,1\}^K$ be the non-adaptive CC-feasible optimum
at level $\delta$, and let $B' = B - z\,\sigma_c\sqrt{T}$ with
$z = \Phi^{-1}(1-\delta/3)$. If \CBUC~only plays arms with
\[
  \sum_{s < t} c_{i_s}
   \;+\; \hat{c}_{i}(t) + \beta\,r_i(t)
   \;\le\; B',
\]
then with probability at least $1 - \delta/3$ the total realised cost
$\sum_{s \le T} c_{i_s}$ does not exceed $B$. \emph{Proof.} The
$z\sigma_c\sqrt{T}$ shrinkage is exactly the Gaussian quantile at
level $\delta/3$ for the sum of $T$ i.i.d.\ sub-Gaussian summands with
variance proxy $\sigma_c^2$; by the sub-Gaussian tail bound,
$\Prob[\sum c_i > \Ex[\sum c_i] + z\sigma_c\sqrt{T}] \le \delta/3$. The
per-round UCB check ensures $\Ex[\sum c_i] \le B'$ on the event of L1,
which itself has probability at least $1 - \delta/3$. A union bound
over the two events gives total failure probability at most
$2\delta/3$. An analogous argument for latency, combined by union bound
with the previous two events, gives a total violation probability
bounded by $\delta$. \qed

\paragraph{L3 (Lagrangian reduction to unconstrained \BwK).}
Let $\mathcal{L}_i(\lambda) = \mu_i - \lambda_B c_i - \lambda_D \ell_i$
with $\lambda = (\lambda_B, \lambda_D) \ge 0$. By Lagrangian duality on
the offline chance-constrained knapsack, there exists a
$\lambda^\star \ge 0$ such that the optimal non-adaptive CC-feasible
value equals $\sup_i \mathcal{L}_i(\lambda^\star) \cdot T$
(up to a rounding gap $O(M)$ from the 0/1 structure). \emph{Proof
sketch.} This is a restatement of LP-duality for 0/1 knapsack with
resource constraints, as used by~\cite{agrawal2014bwk,bwk2018};
the only novelty is applying it to the CC-feasible primal rather than
the mean-feasible primal, which is legitimate because L2 makes the CC
primal equivalent to a mean-feasible primal with modified budget $B'$.
The utility-per-cost ratio that appears in the UCB score
(Eq.~\ref{eq:ucb}) is exactly the per-unit-resource Lagrangian at the
primal optimum of $\lambda^\star_B$ with $\lambda^\star_D$ set to
zero for the binding-budget regime; the non-concavity of the ratio is
sidestepped because we compete only against the non-adaptive oracle,
which itself plays the same ratio. \qed

\paragraph{L4 (\BwK-style regret on the modified instance).}
On the event of L1 $\cap$ L2 (which holds with probability at least
$1-\delta$), the instance reduces to a two-resource \BwK~problem with
modified budget $B'$ and the reduced reward $\mathcal{L}_i(\lambda^\star)$.
Applying~\cite[Thm.~4.1]{bwk2018} to this instance — the theorem
applies directly because the concentration inputs are exactly those
provided by L1 — gives, for some instance-independent $c_1 > 0$,
\[
  \mathsf{REG}_{\mathcal{L}}(T)
   \;\le\; c_1 \sqrt{K T \log T}
        \;+\; \Oh(1).
\]
Translating back from $\mathcal{L}$-regret to $U$-regret adds a
correction proportional to $|B - B'| = z\sigma_c\sqrt{T}$ on the budget
axis and $z\sigma_\ell\sqrt{T}$ on the latency axis: when the
shrunk-budget oracle cannot execute its full plan, the lost utility is
bounded by the Lipschitz constant of $U$ in total resource (at most
$M/c_{\min}$) times the missed budget. Combining,
\[
  \mathsf{REG}(T)
   \;\le\; c_1 \sqrt{K T \log T}
        \;+\; \tfrac{M}{c_{\min}} \cdot z(\sigma_c + \sigma_\ell)\sqrt{T}
        \;+\; c_3.
\]
This is the bound of Theorem~\ref{thm:regret} with
$C_1 = c_1$, $C_2 = M(\sigma_c + \sigma_\ell)/c_{\min}$, $C_3 = c_3$.
The overall failure probability is $\le \delta + KT^{-2}$ (from L1's
union bound), as stated. \qed

\paragraph{Remarks on the nontrivial steps.}
Three steps merit explicit comment. (i) The utility-per-cost ratio in
Eq.~\eqref{eq:ucb} is handled by L3's duality reduction — the ratio
itself is not bounded in regret, but the Lagrangian-linearised score
$\mathcal{L}_i$ is, and L3 shows these coincide at the oracle. (ii) The
two-resource coupling is handled by running the L1--L2 argument
\emph{independently} for cost and latency with confidence level
$\delta/2$ each, then union-bounding; because the UCB check in
\CBUC~uses the max of the two violations, no further coupling term
appears. (iii) The chance-constraint tightening is a finite,
$\sqrt{T}$-scale correction (L2), which is strictly smaller than the
$\sqrt{KT\log T}$ \BwK~term as soon as $K \ge (z\sigma/c_{\min})^2/\log T$;
for all sweep cells studied, this holds by a wide margin, so the
dominant $\tOh(\sqrt{KT})$ rate is preserved.

\paragraph{Proof of Corollary~\ref{cor:cvar} (heavy-tailed cost).}
Replace L2's Gaussian quantile with the empirical-CVaR estimator
$\widehat{\mathrm{CVaR}}_\delta(c_i)$, which under (A1') is consistent
at rate $\Oh(\sqrt{\log t / N_i})$. The per-round feasibility check
becomes $\sum_{s<t} c_{i_s} + \widehat{\mathrm{CVaR}}_\delta(c_i) \le B$.
Repeating L2 with the extreme-order-statistic concentration rate (in
place of the sub-Gaussian quantile) gives the same feasibility
guarantee at the price of a $\sqrt{\log T}$ factor in the correction
term, so the $C_2\,z\,\sigma\sqrt{T}$ term in Theorem~\ref{thm:regret}
becomes $C'_2\,\sqrt{T \log T}$. Steps L1, L3, L4 are unchanged.
\qed

\section{Ablation Studies}
\label{app:ablations}

We report three ablation studies that complement the main evaluation
(\S\ref{sec:evaluation}). All experiments use the same simulator
(\texttt{cbslao\_sim.py}) and are fully reproducible from
\texttt{code/ablations.py}.

% ── A1: α/β/γ sensitivity ──────────────────────────────────────
\subsection{A1: Sensitivity to $\alpha$, $\beta$, $\gamma$}
\label{app:abl-abg}

We perform a one-at-a-time sweep of the three CBUC confidence-width
parameters around their defaults ($\alpha\!=\!1.0$, $\beta\!=\!1.5$,
$\gamma\!=\!1.5$) with $K\!=\!50$, $\rho\!=\!0.5$, tail index~$2.0$,
sparsity~$0.2$, and $n\!=\!40$ seeds per setting.

\begin{table}[t]
\centering
\caption{CBUC sensitivity to $\alpha$, $\beta$, $\gamma$.
Budget-overrun rate remains \textbf{0.000} across the entire grid;
utility varies within a modest range.}
\label{tab:abl-abg}
\begin{tabular}{llcc}
\toprule
Param & Value & Overrun & Utility \\
\midrule
$\alpha$ & 0.25 & 0.000 & 0.792 \\
$\alpha$ & 0.50 & 0.000 & 0.652 \\
$\alpha$ & 1.00 & 0.000 & 0.645 \\
$\alpha$ & 1.50 & 0.000 & 0.670 \\
$\alpha$ & 2.00 & 0.000 & 0.695 \\
$\alpha$ & 3.00 & 0.000 & 0.620 \\
\midrule
$\beta$  & 0.50 & 0.000 & 0.733 \\
$\beta$  & 1.00 & 0.000 & 0.605 \\
$\beta$  & 1.50 & 0.000 & 0.711 \\
$\beta$  & 2.00 & 0.000 & 0.687 \\
$\beta$  & 2.50 & 0.000 & 0.852 \\
$\beta$  & 3.00 & 0.000 & 0.801 \\
\midrule
$\gamma$ & 0.50 & 0.000 & 0.519 \\
$\gamma$ & 1.00 & 0.000 & 0.592 \\
$\gamma$ & 1.50 & 0.000 & 0.684 \\
$\gamma$ & 2.00 & 0.000 & 0.721 \\
$\gamma$ & 2.50 & 0.000 & 0.681 \\
$\gamma$ & 3.00 & 0.000 & 0.762 \\
\bottomrule
\end{tabular}
\end{table}

\paragraph{Observation.}
Table~\ref{tab:abl-abg} shows that the budget-overrun rate is
\emph{identically zero} across the entire parameter grid, confirming
that \CBUC's safety guarantee is insensitive to reasonable
perturbations of $\alpha$, $\beta$, $\gamma$. Utility varies within the
range $[0.519, 0.852]$ but never at the expense of constraint
satisfaction. Figure~\ref{fig:abl-abg} visualises the trade-off: higher
$\gamma$ values (more conservative cost lower bounds) tend to improve
utility by steering the policy toward reliably cheap, high-value tools.

\begin{figure}[t]
\centering
\includegraphics[width=\linewidth]{ablation_abg.png}
\caption{CBUC sensitivity to $\alpha$ (utility bonus), $\beta$
(cost/deadline margin), and $\gamma$ (cost lower bound). Left axis:
budget-overrun rate (identically zero). Right axis: utility achieved.
$K\!=\!50$, $\rho\!=\!0.5$, $n\!=\!40$ seeds.}
\label{fig:abl-abg}
\end{figure}


% ── A2: Utility-dense distractors ──────────────────────────────
\subsection{A2: Behaviour Under Utility-Dense Distractors}
\label{app:abl-distract}

This study varies the fraction of tools with non-zero marginal utility
(\emph{utility sparsity}) from $0.1$ (90\% distractors) to $1.0$ (no
distractors), with $K\!=\!50$, $\rho\!=\!0.5$, tail~$2.0$, $n\!=\!30$
seeds per cell.

\begin{table}[t]
\centering
\caption{Effect of distractor density on four policies.
\CBUC~variants maintain zero overrun across all sparsity levels;
baselines degrade as distractors increase.}
\label{tab:abl-distract}
\begin{tabular}{llcc}
\toprule
Policy & Sparsity & Overrun & Utility \\
\midrule
\ReAct-capped     & 0.10 & 0.200 & 0.754 \\
\ReAct-capped     & 0.25 & 0.133 & 0.867 \\
\ReAct-capped     & 0.50 & 0.033 & 0.983 \\
\ReAct-capped     & 0.75 & 0.133 & 0.996 \\
\ReAct-capped     & 1.00 & 0.067 & 1.000 \\
\midrule
\BwK-vanilla      & 0.10 & 0.033 & 0.755 \\
\BwK-vanilla      & 0.25 & 0.133 & 0.917 \\
\BwK-vanilla      & 0.50 & 0.067 & 0.981 \\
\BwK-vanilla      & 0.75 & 0.200 & 0.998 \\
\BwK-vanilla      & 1.00 & 0.067 & 1.000 \\
\midrule
\CBUC-aggr ($z\!=\!0$) & 0.10 & 0.000 & 0.387 \\
\CBUC-aggr ($z\!=\!0$) & 0.25 & 0.000 & 0.775 \\
\CBUC-aggr ($z\!=\!0$) & 0.50 & 0.000 & 0.979 \\
\CBUC-aggr ($z\!=\!0$) & 0.75 & 0.000 & 0.998 \\
\CBUC-aggr ($z\!=\!0$) & 1.00 & 0.000 & 0.998 \\
\midrule
\textbf{\CBUC~($z\!=\!1.28$)} & 0.10 & \textbf{0.000} & 0.437 \\
\textbf{\CBUC~($z\!=\!1.28$)} & 0.25 & \textbf{0.000} & 0.791 \\
\textbf{\CBUC~($z\!=\!1.28$)} & 0.50 & \textbf{0.000} & 0.940 \\
\textbf{\CBUC~($z\!=\!1.28$)} & 0.75 & \textbf{0.000} & 0.996 \\
\textbf{\CBUC~($z\!=\!1.28$)} & 1.00 & \textbf{0.000} & 0.999 \\
\bottomrule
\end{tabular}
\end{table}

\paragraph{Observation.}
Table~\ref{tab:abl-distract} reveals a clear separation: both \CBUC\
variants achieve \emph{zero} budget overrun regardless of distractor
density, whereas \ReAct-capped and \BwK-vanilla reach overrun rates of
up to $20\%$ when 90\% of tools are distractors (sparsity~$=\!0.1$).
The utility gap between \CBUC\ and the baselines is largest when
distractors dominate: at sparsity $0.1$, \CBUC\ achieves $0.437$
vs.\ $0.754$ for \ReAct-capped, but the baselines violate the budget to
obtain that utility. As sparsity increases toward $1.0$, the gap
closes because all tools carry useful payoff and budget pressure
subsides (Figure~\ref{fig:abl-distract}).

\begin{figure}[t]
\centering
\includegraphics[width=\linewidth]{ablation_distractors.png}
\caption{Effect of distractor density. Left: utility achieved.
Right: budget-overrun rate. \CBUC~variants maintain zero overrun at all
densities; baselines degrade under sparse utility. $K\!=\!50$,
$\rho\!=\!0.5$, $n\!=\!30$.}
\label{fig:abl-distract}
\end{figure}


% ── A3: Calibrated benchmark-trace replay ──────────────────────
\subsection{A3: Calibrated Benchmark-Trace Replay}
\label{app:abl-replay}

To gauge \CBUC's behaviour under workloads resembling public
tool-calling benchmarks, we calibrate the simulator to summary
statistics from three benchmark families:
\begin{itemize}
  \item \textbf{ToolBench-like} ($K\!=\!100$, 6~tools/task, heavy tail
    $\alpha\!=\!1.8$, sparse utility $s\!=\!0.08$): large API pool,
    heavy-tailed costs, very sparse relevance.
  \item \textbf{$\tau$-bench-like} ($K\!=\!15$, 5~tools/task,
    $\alpha\!=\!2.5$, $s\!=\!0.4$): small pool, moderate tail, denser
    relevance.
  \item \textbf{BFCL-like} ($K\!=\!40$, 2~tools/task, $\alpha\!=\!2.2$,
    $s\!=\!0.15$): mid-size pool, single-function-call-style tasks.
\end{itemize}
Budget is calibrated to tools-per-task $\times$ mean-cost $\div$
budget-multiplier. $n\!=\!50$ seeds per profile.

\emph{Caveat}: these are parameterised simulations, \textbf{not} live
trace replays. Network calls are not issued; only summary statistics
from published benchmark documentation are used. See~\S\ref{sec:threats}.

\begin{table}[t]
\centering
\caption{Calibrated benchmark-trace replay. \CBUC~achieves $\le 2\%$
overrun across all three profiles; baselines reach $44$--$100\%$.
$n\!=\!50$ seeds per profile.}
\label{tab:abl-replay}
\begin{tabular}{llccc}
\toprule
Profile & Policy & Overrun & Utility & Regret \\
\midrule
\multirow{4}{*}{ToolBench}
  & \ReAct-capped     & 0.620 & 0.646 & $+0.324$ \\
  & \BwK-vanilla      & 0.440 & 0.653 & $+0.318$ \\
  & \CBUC-aggr        & \textbf{0.000} & 0.427 & $+0.543$ \\
  & \textbf{\CBUC}    & \textbf{0.000} & 0.469 & $+0.501$ \\
\midrule
\multirow{4}{*}{$\tau$-bench}
  & \ReAct-capped     & 1.000 & 0.802 & $+0.135$ \\
  & \BwK-vanilla      & 0.560 & 0.810 & $+0.127$ \\
  & \CBUC-aggr        & \textbf{0.000} & 0.738 & $+0.199$ \\
  & \textbf{\CBUC}    & \textbf{0.000} & 0.728 & $+0.208$ \\
\midrule
\multirow{4}{*}{BFCL}
  & \ReAct-capped     & 1.000 & 0.387 & $+0.491$ \\
  & \BwK-vanilla      & 0.580 & 0.374 & $+0.504$ \\
  & \CBUC-aggr        & 0.040 & 0.518 & $+0.360$ \\
  & \textbf{\CBUC}    & \textbf{0.020} & 0.492 & $+0.386$ \\
\bottomrule
\end{tabular}
\end{table}

\paragraph{Observation.}
Table~\ref{tab:abl-replay} confirms that baseline policies fail
catastrophically on all three calibration profiles: \ReAct-capped
overruns the budget in $62$--$100\%$ of runs, and \BwK-vanilla in
$44$--$58\%$. In contrast, \CBUC~achieves $\le 2\%$ overrun across
every profile. On the ToolBench-like workload --- the most challenging
profile due to its large pool, heavy tail, and extreme sparsity ---
\CBUC\ still maintains zero overrun at the cost of a moderate utility
reduction ($0.469$ vs.\ $0.653$ for \BwK-vanilla). On the denser
$\tau$-bench profile, the utility gap narrows to $0.728$ vs.\ $0.810$,
and the baselines overrun the budget more than half the time.
Figure~\ref{fig:abl-replay} summarises the comparison.

\begin{figure}[t]
\centering
\includegraphics[width=\linewidth]{ablation_replay.png}
\caption{Calibrated benchmark replay (parameterised to published
summary statistics; not a live-trace replay). Left: budget-overrun
rate. Right: utility achieved. \CBUC~achieves near-zero overrun across
all profiles. $n\!=\!50$.}
\label{fig:abl-replay}
\end{figure}
